% ============================================================================
% GÉNÉRÉ par scripts/exemples-guide.ts — NE PAS ÉDITER À LA MAIN.
% Régénérer : npm run exemples (chaque chiffre est vérifié contre le moteur).
% ============================================================================
\pexemples Les paramètres sont des moyennes annuelles des quatre barèmes
trimestriels (d'où les décimales). Le SRG et son supplément se calculent par
\emph{tranches} de revenu (\num{24}, \num{48} ou
\num{96}~\$) — le revenu de pension est arrondi à la tranche
avant d'appliquer le taux de réduction.

\medskip\noindent\textbf{M10 --- retraité vivant seul, 70~ans, revenu de pension privée de \D{20000}.}
Pension de base : \num{8791.14} (le revenu est loin du seuil de
récupération de \num{93454}). SRG (taux de \pc{50},
tranches de \num{24}~\$) :
\begin{align*}
\text{revenu arrondi} &= \Bigl\lfloor \tfrac{\num{20000}}{\num{24}} \Bigr\rfloor \times \num{24} = \num{19992} \\
SRG &= \pos{\num{11096.85} - \num{19992} \times \pc{50}} = \num{1100.85}
\end{align*}
Le supplément complémentaire est éteint à ce revenu (\num{0}).
Total : \num{8791.14} $+$ \num{1100.85} $=$ \D{9891.99}.

\medskip\noindent\textbf{M13 --- couple de retraités, 68 et 66~ans, pensions privées de \D{12000} et \D{6000}.}
Couple dont les deux conjoints ont 65~ans et plus, revenu combiné de
\num{18000}. Chacun reçoit la pension de base (\num{8791.14}) et le SRG au
barème « couple » (maximum \num{7327.65}, taux \pc{25},
tranches de \num{48}~\$) :
\begin{align*}
SRG &= \pos{\num{7327.65} - \num{18000} \times \pc{25}} = \num{2827.65} \quad (\text{par conjoint}) \\
\text{supplément du couple} &= \pos{\num{1152.60} - \num{13920} \times \pc{25}} = \num{0} \\
\text{total} &= 2 \times \num{8791.14} + 2 \times \num{2827.65} + \num{0} = \num{23237.58}
\end{align*}
Sécurité de la vieillesse du ménage : \D{23237.58}.

\medskip\noindent\textbf{M11 --- couple de retraités, 72 et 70~ans, pensions privées de \D{30000} et \D{10000}.}
Même structure, revenu combiné de \num{40000} : la réduction du SRG
(\num{39984} $\times$ \pc{25} $=$
\num{9996}) dépasse le maximum de
\num{7327.65} : SRG \emph{éteint}. Reste la pension de base de chacun :
$2 \times \num{8791.14} = \num{17582.28}$, soit \D{17582.28}.

\medskip\noindent\textbf{M12 --- retraité vivant seul, 76~ans, revenu de pension privée de \D{110000}.}
75~ans et plus : pension majorée (\num{8791.14} $+$ \num{879.15}
$=$ \num{9670.29}). Le revenu déclenche l'\emph{impôt de récupération} :
\begin{align*}
\text{récupération} &= \pos{\num{110000} + \num{9670.29} - \num{93454}} \times \pc{15} = \num{3932.44} \\
PSV &= \num{9670.29} - \num{3932.44} = \num{5737.85}
\end{align*}
Pension nette : \D{5737.85} (la pension s'éteint complètement vers
\D{148252.31} de revenu).
